% ---------- Premises: every problem type in the unit, what is given, what is asked ----------
\bpart[START2]{START}{Premises: every problem type, what you are given and what you must find}{1}
{\small Read the question, find its row, write the \emph{given} and the \emph{find} at the top of your answer (it earns method marks), then use the core relation. Symbols: $\mathbf q=(q_1,\dots,q_n)$ joint variables ($q_i=\theta_i$ for a revolute joint in rad, $q_i=d_i$ for a prismatic joint in m); \{0\} or \{S\} base/station frame; \{E\} or \{T\} end-effector/tool frame; ${}^A\T_B$ pose of \{B\} expressed in \{A\}; subscript $d$ = desired.\par}
\vspace{2pt}
{\small\renewcommand{\arraystretch}{1.12}
\begin{tabularx}{\linewidth}{@{}>{\raggedright\arraybackslash}p{0.13\linewidth}>{\raggedright\arraybackslash}X>{\raggedright\arraybackslash}X>{\raggedright\arraybackslash}p{0.25\linewidth}>{\raggedright\arraybackslash}p{0.075\linewidth}@{}}
\toprule
\textbf{Problem} & \textbf{Given} & \textbf{Find} & \textbf{Core relation} & \textbf{Go to}\\\midrule
\multicolumn{5}{@{}l}{\sffamily\bfseries\color{seckin}Mechanisms and kinematics (W1–W6)}\\
Mobility (DoF) of a mechanism & a drawing: links $N$ (count the ground), joints $J$, freedom $f_i$ of each joint; planar ($D=3$) or spatial ($D=6$) & mobility $M$ = number of independent inputs needed & $M=D(N-J-1)+\sum_{i=1}^{J}f_i$; extended: $+R_C-R_M$ & W1\,p.\pageref{p:W1}, W4\,p.\pageref{p:W4}\\
Forward kinematics (FK) & joint variables $\mathbf q$; geometry (link lengths, DH table or a drawing) & tool pose ${}^0\T_E=\begin{bmatrix}{}^0\Rm_E&{}^0\mathbf p_E\\\mathbf 0^\top&1\end{bmatrix}$ (position and orientation) & ${}^0\T_E={}^0\T_1(q_1)\,{}^1\T_2(q_2)\cdots{}^{n-1}\T_n(q_n)\,{}^n\T_E$ (ETS or DH) & W2\,p.\pageref{p:W2}\\
Transform between frames & frames drawn (or described) relative to each other & ${}^A\T_B$, a point in another frame ${}^A\mathbf P={}^A\T_B\,{}^B\mathbf P$, or the inverse ${}^B\T_A$ & ${}^A\T_C={}^A\T_B\,{}^B\T_C$; ${}^B\T_A=\begin{bmatrix}\Rm^\top&-\Rm^\top\mathbf t\\\mathbf 0^\top&1\end{bmatrix}$ & W2\,p.\pageref{p:W2}\\
Inverse kinematics (IK) & desired tool pose ${}^0\T_E$, or only position $(x,y,z)$ / planar $(x,y,\phi)$ & \emph{all} joint vectors $\mathbf q$ that reach it (0, 1, 2, … solutions; say which branch) & solve $\mathrm{FK}(\mathbf q)={}^0\T_E$: geometric (cosine rule, $\atantwo$) or algebraic & W3\,p.\pageref{p:W3}\\
Workspace & link lengths, joint types (and limits) & reachable workspace (every position the tool can reach) / dexterous (reach with every orientation); sketch it & distance range from the shoulder $|L_1-L_2|\dots L_1+L_2$; sweep about the base axis & W3\,p.\pageref{p:W3}\\
Parallel (closed-chain) IK & desired output pose of the platform / end link (e.g. four-bar output angle) & actuated (and passive) joint angles, every assembly mode & loop-closure: the vector sum round each loop $=\mathbf 0$ & W4\,p.\pageref{p:W4}\\
Parallel FK & actuated joint values & platform pose (usually several solutions) & same loop-closure equations, solved the other way & W4\,p.\pageref{p:W4}\\
Trajectory & start, goal (and via) values $\theta_0,\theta_f$ ($\theta_v$), duration $t_f$, boundary velocities (accelerations), or limits $\dot\theta_{\max}$, $\ddot\theta$ & $\theta(t),\dot\theta(t),\ddot\theta(t)$: cubic coefficients $a_0..a_3$, or LSPB blend time $t_b$ & cubic: 4 boundary conditions → 4 coefficients; LSPB: $s=ut+\tfrac12at^2$ per segment & W5\,p.\pageref{p:W5}\\
Velocity kinematics & $\mathbf q$, joint rates $\dot{\mathbf q}$ & tool velocity $\boldsymbol\nu=(\boldsymbol\upsilon,\boldsymbol\omega)$ (linear, angular) & $\boldsymbol\nu=\mathbf J(\mathbf q)\,\dot{\mathbf q}$ & W6\,p.\pageref{p:W6}\\
Inverse velocity (resolved rate) & desired tool velocity $\boldsymbol\nu$ (e.g. trace a line/circle), $\mathbf q$ & joint rates $\dot{\mathbf q}$ & $\dot{\mathbf q}=\mathbf J^{-1}(\mathbf q)\,\boldsymbol\nu$ (square, non-singular $\mathbf J$) & W6\,p.\pageref{p:W6}\\
Singularities & $\mathbf J(\mathbf q)$ & poses where the arm loses a direction of motion & $\det\mathbf J(\mathbf q)=0$ & W6\,p.\pageref{p:W6}\\
Statics & wrench at the tool $\mathbf F_E$ (force N, moment N\,m), $\mathbf q$ & joint torques/forces $\boldsymbol\tau$ that hold it & $\boldsymbol\tau=\mathbf J^\top(\mathbf q)\,\mathbf F_E$ & W6\,p.\pageref{p:W6}\\[2pt]
\multicolumn{5}{@{}l}{\sffamily\bfseries\color{secdyn}Dynamics and control (W7–W9)}\\
Equations of motion (EoM) & masses $m_i$, centre-of-mass positions, inertias $I_i$, geometry, gravity $g$ & $\mathbf Q=\mathbf M(\mathbf q)\ddot{\mathbf q}+\mathbf C(\mathbf q,\dot{\mathbf q})+\mathbf g(\mathbf q)$ in full & Lagrange: $\mathcal L=K-P$, $Q_i=\frac{d}{dt}\frac{\partial\mathcal L}{\partial\dot q_i}-\frac{\partial\mathcal L}{\partial q_i}$; or Newton–Euler & W7\,p.\pageref{p:W7}\\
Inverse dynamics & motion $\mathbf q,\dot{\mathbf q},\ddot{\mathbf q}$ & actuator torques/forces $\mathbf Q$ & substitute into the EoM & W7\,p.\pageref{p:W7}, W9\,p.\pageref{p:W9}\\
Forward dynamics & torques $\mathbf Q$, state $\mathbf q,\dot{\mathbf q}$ & accelerations $\ddot{\mathbf q}$ (then integrate) & $\ddot{\mathbf q}=\mathbf M^{-1}(\mathbf Q-\mathbf C-\mathbf g)$ & W7\,p.\pageref{p:W7}\\
Single-joint control (servo) & plant $m\ddot x+b\dot x+kx=f$, specification (critically damped, stiffness, disturbance) & gains $k_p,k_v$ ($k_i$), closed-loop response, steady-state error & $f=\alpha f'+\beta$, $f'=\ddot x_d+k_v\dot E+k_pE$; critical: $k_v=2\sqrt{k_p}$ & W8\,p.\pageref{p:W8}\\
Model-based (computed-torque) control & robot model $\mathbf M,\mathbf C,\mathbf g$; desired trajectory $\mathbf q_d(t)$ & control law $\mathbf Q$; error dynamics & $\mathbf Q=\boldsymbol\alpha\acute{\mathbf Q}+\boldsymbol\beta$, $\boldsymbol\alpha=\mathbf M(\mathbf q)$, $\boldsymbol\beta=\mathbf C+\mathbf g$, $\acute{\mathbf Q}=\ddot{\mathbf q}_d+\mathbf K_v\dot{\mathbf E}+\mathbf K_p\mathbf E$ & W9\,p.\pageref{p:W9}\\
Cartesian / force control & desired Cartesian path $\mathbf x_d(t)$ or contact force $\mathbf f_d$ & joint torques; which directions are position- vs force-controlled & $\mathbf Q=\mathbf J^\top\mathcal F$; error $\mathbf E=\mathbf x_d-\mathbf x$ or $\mathbf f_d-\mathbf f_e$ & W9\,p.\pageref{p:W9}\\[2pt]
\multicolumn{5}{@{}l}{\sffamily\bfseries\color{secmob}Wheeled robots (W10–W11)}\\
Wheeled FK & wheel speeds $\dot\varphi_i$ (and steering angles $\beta_i$), wheel layout & chassis velocity $\zeta=(u,v,\dot\theta)$ in \{C\}; world velocity $(\dot x,\dot y,\dot\theta)$ in \{O\}; ICR & invert the stacked wheel equations; $\nu=\mathbf J(\theta)\,\zeta$ & W10\,p.\pageref{p:W10}\\
Wheeled IK & desired chassis (or world) velocity, e.g. drive a circle or a square & wheel speeds $\dot\varphi_i$ (and steering) & each wheel's rolling equation, one row per wheel & W10\,p.\pageref{p:W10}\\
Wheeled-robot dynamics & mass $m$, inertia ${}^GI$, wheel forces, COM offset & chassis accelerations; the forces/torques needed & Newton–Euler in \{C\}: $\mathbf F=m\,\mathbf a_G$, $N_z={}^GI\ddot\theta$ & W11\,p.\pageref{p:W11}\\[2pt]
\multicolumn{5}{@{}l}{\sffamily\bfseries\color{secplan}Planning (W12)}\\
Graph search & graph/grid, edge costs, start, goal (heuristic $h$) & order of expansion, path, its cost; is it optimal? & BFS (queue), DFS (stack), Dijkstra (cost $g$), A* ($f=g+h$, $h$ admissible) & W12\,p.\pageref{p:W12}\\
Sampling-based planning & C-space, obstacles, start, goal & a collision-free path (not optimal) & PRM (roadmap, then search); RRT (grow a tree towards samples) & W12\,p.\pageref{p:W12}\\
\bottomrule
\end{tabularx}}
